% !TEX root = ../main.tex
\chapter{Resultats de qualitat i compressió}\label{ch:resultats}
\fancyhead[RE]{\slshape \textsf{Capítol \thechapter.\ Resultats}}
\fancyhead[LO]{\slshape \textsf{\nouppercase \rightmark}}

\section{Lectura del capítol}\label{sec:result_reading}

Els resultats segueixen el mateix ordre que la metodologia. Primer es presenta
el full C-only, que descriu cobertura i serveix de cribratge. Després es mostra
CSMR, la ruta de referència de SORTENY basada en TensorFlow Compression i
l'única font de bitrate final. Finalment s'analitzen resultats per preset,
modes experimentals, \code{preserve-roi}, evidència visual i regles d'ús.

Les comparacions principals utilitzen \code{q204}. Quan apareix
\code{adaptive\_s8}, s'identifica com a baseline secundari. Tots els deltes ROI
i fons respecte de \code{q204} s'han recalculat amb la mateixa màscara del
preset.

\section{Validació del full C-only}\label{sec:full_results}

Els 120 retalls es van processar i van produir 4.800 registres principals, sense
fallades reportades. Aquest nombre no significa 4.800 imatges diferents, sinó
les combinacions de 10 presets, 120 retalls i 4 polítiques. Tots els Q-maps
tenien 1.024 bytes i les reconstruccions conservades tenien la mida esperada.

\subsection{Resultats de cobertura intermèdia}

La Taula~\ref{tab:c_only_mid} resumeix \code{focus\_bgq128} en els casos
de cobertura intermèdia, és a dir, casos on la ROI ocupa entre el 10\% i el
70\% dels blocs. En aquesta franja hi ha prou regió protegida i prou fons per
mesurar redistribució. Totes les files comparen \code{focus\_bgq128} amb
\code{q204} del mateix retall i la mateixa màscara ROI/fons.

Les magnituds de la taula són:

\begin{itemize}[leftmargin=*]
  \item $\Delta$ROI =
  $\mathrm{PSNR}_{ROI}(\mathrm{focus})-\mathrm{PSNR}_{ROI}(\mathrm{q204})$;
  \item $\Delta$fons =
  $\mathrm{PSNR}_{fons}(\mathrm{focus})-\mathrm{PSNR}_{fons}(\mathrm{q204})$;
  \item $\Delta\bar Q =
  \bar Q(\mathrm{focus})-\bar Q(\mathrm{q204})$;
  \item $\Delta H =
  H_{latent}(\mathrm{focus})-H_{latent}(\mathrm{q204})$.
\end{itemize}

$\Delta H$ és un proxy C-only en bits per símbol latent. No és bitrate real; el
bitrate real es mesura després amb CSMR. \code{adaptive\_s8} es conserva com a
control tècnic perquè redistribueix Q amb mitjana propera a 204, però no és el
baseline principal d'aquesta taula.

\begin{table}[H]
\centering
\scriptsize
\begin{tabular}{lrrrrrr}
\toprule
Preset & Registres & ROI \% & $\Delta$ROI & $\Delta$fons & $\Delta\bar Q$ &
$\Delta H$ \\
\midrule
\code{vegetation} & 29 & 33,6 & +0,257 & -2,069 & -44,75 & -0,206 \\
\code{chlorophyll} & 5 & 18,9 & +0,168 & -2,088 & -58,47 & -0,269 \\
\code{cloud\_avoid} & 13 & 43,1 & +0,294 & -2,089 & -36,16 & -0,154 \\
\code{local\_contrast} & 63 & 35,9 & +0,195 & -2,267 & -42,91 & -0,185 \\
\code{vegetation\_green} & 25 & 36,7 & +0,270 & -2,064 & -41,92 & -0,193 \\
\code{clouds} & 23 & 26,9 & +0,160 & -2,115 & -51,33 & -0,234 \\
\code{low\_ndvi} & 17 & 28,9 & +0,173 & -2,080 & -49,67 & -0,236 \\
\bottomrule
\end{tabular}
\caption{Cribratge C-only per a \code{focus\_bgq128} en casos de cobertura
intermèdia. \emph{Registres} és el nombre de combinacions retall--preset que
compleixen el filtre. El baseline és \code{q204}; els deltes de PSNR són en dB
i $\Delta H$ és en bits per símbol latent.}
\label{tab:c_only_mid}
\end{table}

La taula mostra un patró consistent: la ROI guanya aproximadament entre
0,16 i 0,29~dB, mentre que el fons perd al voltant de 2~dB i la Q mitjana
baixa. Això valida
el mecanisme de redistribució i justifica passar casos a CSMR. No demostra
encara reducció del fitxer codificat. El ranking detallat per proxies queda
documentat a l'Apèndix~\ref{app:evidence}; en aquest capítol només s'utilitza
com a criteri de cribratge, no com a resultat final.

\subsection{Cobertura com a resultat}

\code{local\_contrast} produeix 63 casos mid-roi i, per tant, és el preset amb
més cobertura útil en aquest dataset. \code{vegetation} en produeix 29 i
\code{vegetation\_green}, 25. \code{chlorophyll} només en produeix cinc. La
quantitat de registres no mesura qualitat del preset; mesura quantes escenes
disponibles generen una separació ROI/fons dins l'interval 10--70\%.

\code{clouds} apareix amb resultats C-only, però es manté experimental perquè
no disposa de B11. \code{cloud\_avoid} té la mateixa limitació espectral, però
la seva interpretació operativa --protegir zones no ennuvolades o útils-- és
coherent amb l'objectiu de no gastar bits en núvols probables.

\section{Precisió de \texorpdfstring{\code{global\_target}}{global target}}\label{sec:global_target_results}

\code{global\_target} respon una pregunta diferent de les polítiques amb ROI:
no intenta redistribuir qualitat dins la imatge, sinó convertir un objectiu
global de PSNR o MSE en una Q uniforme. Per això s'analitza com a validació de
predicció, no com a competidor de \code{focus\_bgq128}.

La Figura~\ref{fig:global_target_curve} mostra la corba mesurada Q--PSNR per a
la calibració global canonical en l'escala actual \code{lambda005}. Aquí
\emph{rang calibrat} vol dir l'interval de PSNR observat quan es fa variar Q en
la zona principal Q128--Q255. Si un objectiu cau dins aquest interval, es pot
interpolar una Q constant. Si l'objectiu queda per sobre de la corba mesurada,
el mode no pot assolir-lo amb una Q uniforme i satura a Q255.

\begin{figure}[H]
\centering
\includegraphics[width=0.90\textwidth]{global_target_lambda005_curves.png}
\caption{Corba Q--PSNR canonical de \code{global\_target} en la recalibració
\code{lambda005}. La línia horitzontal marca un objectiu de 76,8~dB. Com que la
corba canonical arriba només a 74,52~dB amb Q255, aquest target queda fora del
rang mesurat i la decisió satura.}
\label{fig:global_target_curve}
\end{figure}

La Taula~\ref{tab:global_target_precision} resumeix aquesta lectura per a
l'objectiu 76,8~dB. La idea important és que \code{global\_target} no força
qualsevol PSNR demanada: només pot seleccionar una Q dins del rang mesurat per
la calibració disponible.

\begin{table}[H]
\centering
\small
\begin{tabular}{lll}
\toprule
Magnitud & Valor & Lectura \\
\midrule
Rang PSNR Q128--Q255 & 71,64--74,52~dB & Rang calibrat canonical \\
\code{q204} & 73,62~dB & Referència uniforme dins del rang \\
Q255 & 74,52~dB & Màxim mesurat en aquesta calibració \\
Objectiu 76,8~dB & Fora de rang & Satura a Q255 \\
\bottomrule
\end{tabular}
\caption{Resultats de l'auditoria de \code{global\_target}. La població són
els punts Q128--Q255 de la calibració canonical \code{lambda005}. No es compara
contra \code{q204}: valida la conversió objectiu--Q i els seus límits.}
\label{tab:global_target_precision}
\end{table}

La conclusió és útil per a la resta de la memòria: la predicció global és
interpretable quan l'objectiu cau dins el rang calibrat, però no pot inventar
qualitat quan el target queda fora de la corba mesurada. Aquesta és la mateixa
raó per la qual les polítiques espacials es validen després amb reconstrucció,
PSNR ROI/fons i bitrate \code{.tfci}: la Q-map és una decisió de qualitat, però
el resultat final s'ha de mesurar en una execució real.

L'abast d'aquesta secció és deliberadament limitat. No és una estadística
multiimatge de \code{global\_target}, sinó una comprovació del mecanisme amb la
calibració global canònica: mostra com es transforma un objectiu en una Q
uniforme i què passa quan l'objectiu queda fora del rang mesurat. Una avaluació
estadística de precisió global requeriria repetir aquesta auditoria sobre més
retalls i objectius.

\section{Validació CSMR}\label{sec:csmr_results}

Els resultats d'aquesta secció no són prediccions de la taula de calibració. Són
execucions reals del codificador de referència de SORTENY, implementat amb
TensorFlow Compression, amb els Q-maps generats prèviament. La taula de
calibració ajuda a decidir quins valors Q escriure; CSMR mesura què passa quan
aquests valors es transformen en l'array local de qualitat i es codifiquen els
latents en un fitxer \code{.tfci}.

La mostra principal es va construir a partir del cribratge C-only: fins a cinc
casos de \code{cloud\_avoid}, fins a cinc de \code{local\_contrast}, fins a
quatre de \code{vegetation}, casos addicionals de \code{vegetation\_green} i
\code{chlorophyll}, i almenys un cas límit o negatiu quan era disponible. La
selecció no busca estimar una mitjana universal; busca cobrir situacions on la
redistribució de qualitat és plausible i casos on podria fallar.

\subsection{Correspondència L igual Q}

Els 20 retalls seleccionats generen 60 subcasos: tres polítiques per cada
retall. Tots van completar compressió i descompressió. Els 60 arrays tenien
4.096 bytes i els 60 fitxers \code{.tfci}
van produir reconstruccions amb forma $8\times512\times512$. La comprovació
clau va donar:

\begin{equation}
\code{packed\_l\_matches\_qmap}=60/60.
\end{equation}

Això confirma que l'array local de qualitat aplica per bloc els valors derivats
del Q-map C. No demostra per si sol que el preset sigui bo; demostra que
l'experiment mesura la política que s'ha sol·licitat i no una qualitat global
equivalent. Tampoc valida que la predicció de MSE de la calibració sigui exacta
en cada bloc; aquesta bondat es jutja amb PSNR, ROI/fons i bitrate \code{.tfci}.

\subsection{Comparació global de polítiques}

\begin{table}[H]
\centering
\small
\begin{tabular}{lrrrr}
\toprule
Política & Subcasos & Bitrate (bps) & PSNR global (dB) & $L=Q$ \\
\midrule
\code{q204} & 20 & 4,0602 & 74,272 & 100\% \\
\code{adaptive\_s8} & 20 & 4,0634 & 74,291 & 100\% \\
\code{focus\_bgq128} & 20 & 3,7013 & 72,369 & 100\% \\
\bottomrule
\end{tabular}
\caption{Bitrate real CSMR i PSNR global en els 20 retalls seleccionats. Cada
fila agrega 20 subcasos, un per retall. El bitrate és en bps, bits per mostra
espectral codificada.}
\label{tab:csmr_global}
\end{table}

\begin{figure}[H]
\centering
\includegraphics[width=0.72\textwidth]{csmr_policy_bitrate.png}
\caption{Bitrate mitjà dels tres controls CSMR sobre els mateixos 20 retalls.
La comparació principal és contra \code{q204}; \code{adaptive\_s8} s'inclou com
control de Q variable.}
\label{fig:csmr_policy}
\end{figure}

\code{adaptive\_s8} té pràcticament el mateix bitrate que \code{q204}: 4,0634
enfront de 4,0602~bps. El resultat no implica que el Q-map adaptatiu sigui
idèntic; conté entre Q196 i Q240. La seva mitjana és 203,994 i la redistribució
local no canvia prou l'estadística codificada per reduir la taxa global. Per
tant, \code{adaptive\_s8} és útil com a control de qualitat local, però no aporta
estalvi mesurable en aquesta mostra.

\code{focus\_bgq128} baixa a 3,7013~bps. Calculant el quocient entre mitjanes:

\begin{equation}
100\left(1-\frac{3,7013}{4,0602}\right)=8,84\%.
\end{equation}

El fitxer focus envia el 91,16\% dels bits de \code{q204}, una millora relativa
d'aproximadament $1,097\times$.

\subsection{Qualitat amb la mateixa màscara}

La mitjana de deltes per cas respecte de \code{q204} és:

\begin{equation}
\Delta\mathrm{PSNR}_{ROI}=+0,203\;\mathrm{dB},
\qquad
\Delta\mathrm{PSNR}_{fons}=-2,204\;\mathrm{dB}.
\end{equation}

Respecte d'\code{adaptive\_s8}, els valors són +0,2137 i -2,2296~dB. Les dues
referències condueixen a la mateixa conclusió: el mode conserva o millora la
ROI i obté l'estalvi degradant el fons. El PSNR global baixa perquè el fons
ocupa la major part de mostres; aquesta baixada no s'ha d'interpretar sense
separar regions.

\section{Resultats per preset}\label{sec:preset_results}

\begin{table}[H]
\centering
\scriptsize
\begin{tabular}{lrrrrr}
\toprule
Preset & Subcasos & ROI \% & Estalvi vs \code{q204} & $\Delta$ROI & $\Delta$fons \\
\midrule
\code{local\_contrast} & 5 & 10,6 & 11,68\% & +0,147 & -2,504 \\
\code{cloud\_avoid} & 5 & 20,6 & 11,56\% & +0,280 & -2,158 \\
\code{vegetation} & 4 & 11,6 & 8,56\% & +0,190 & -1,934 \\
\code{vegetation\_green} & 3 & 12,8 & 8,00\% & +0,241 & -2,221 \\
\code{chlorophyll} & 3 & 19,4 & 7,41\% & +0,147 & -2,121 \\
\bottomrule
\end{tabular}
\caption{Combinacions \code{preset + focus\_bgq128} respecte de
\code{q204}. Els deltes de PSNR són en dB i utilitzen la mateixa màscara. Cada
fila agrega els subcasos CSMR indicats.}
\label{tab:preset_csmr}
\end{table}

\begin{figure}[H]
\centering
\includegraphics[width=0.86\textwidth]{csmr_savings_main_presets.png}
\caption{Estalvi mitjà per preset respecte de \code{q204}. El nombre sobre cada
barra és la mida de la mostra CSMR, no el nombre de retalls del full.}
\label{fig:preset_savings}
\end{figure}

\code{local\_contrast} i \code{cloud\_avoid} superen l'11\% de mitjana en les
seves mostres. \code{vegetation}, \code{vegetation\_green} i
\code{chlorophyll} queden entre el 7,4 i el 8,6\%. Tots presenten delta ROI
positiu i fons negatiu. Per tant, els presets no milloren una ROI manual:
milloren la qualitat de la regió que ells mateixos detecten respecte de
comprimir uniformement amb \code{q204}. Aquesta és la contribució operativa
principal.

Les mitjanes per preset no s'han de comparar sense considerar la mostra.
Cada preset es va executar en retalls seleccionats diferents. El 11,68\% de
\code{local\_contrast} no prova que sempre sigui millor que
\code{vegetation}; prova que, en els cinc casos seleccionats, la seva cobertura
permet un compromís més agressiu.

\section{Diagrames taxa--distorsió}\label{sec:rd_results}

\begin{figure}[H]
\centering
\includegraphics[width=0.88\textwidth]{rd_roi_q204_focus.png}
\caption{Diagrama taxa--distorsió dins la ROI. Cada parella unida correspon al
mateix retall i la mateixa màscara. Un desplaçament a l'esquerra i amunt indica
menys bitrate i més PSNR ROI.}
\label{fig:rd_roi}
\end{figure}

La Figura~\ref{fig:rd_roi} mostra que la majoria de punts focus es desplacen a
l'esquerra i lleugerament amunt respecte de \code{q204}. La dispersió vertical
entre retalls és molt més gran que el delta d'una política, perquè les escenes
tenen errors base diferents. Per això les línies per parella són més
informatives que comparar núvols de punts sense correspondència.

\begin{figure}[H]
\centering
\includegraphics[width=0.88\textwidth]{rd_background_q204_focus.png}
\caption{Diagrama taxa--distorsió al fons. El moviment esperat és cap a
l'esquerra i avall: menys bitrate a canvi de més distorsió en la regió no
prioritària.}
\label{fig:rd_background}
\end{figure}

Al fons, Figura~\ref{fig:rd_background}, el desplaçament cap avall és clar.
Aquest resultat evita una interpretació equivocada: focus no millora tota la
imatge mentre comprimeix més. Redistribueix un pressupost de qualitat.

\section{Modes experimentals}\label{sec:experimental_results}

L'auditoria experimental va completar 30 de 30 subcasos i va conservar $L=Q$
en tots. Cada fila de la Taula~\ref{tab:experimental} compara el mode avaluat
contra \code{q204} del mateix retall i la mateixa màscara ROI/fons.

\begin{table}[H]
\centering
\small
\begin{tabular}{lrrrrl}
\toprule
Mode avaluat & Subcasos & ROI \% & Estalvi & $\Delta$ROI & Lectura \\
\midrule
\code{high\_ndvi + focus\_bgq128} & 3 & 20,4 & 6,54\% & +0,165 & Candidat útil \\
\code{low\_ndvi + focus\_bgq128} & 3 & 4,9 & 9,05\% & +0,120 & ROI massa petita \\
\code{dark\_regions + focus\_bgq128} & 3 & 62,4 & 1,99\% & +0,384 & ROI massa ampla \\
\code{water\_body + focus\_bgq128} & 1 & 10,1 & 11,18\% & +0,158 & Una sola mostra \\
\bottomrule
\end{tabular}
\caption{CSMR dels llindars experimentals. Estalvi i PSNR es calculen respecte
de \code{q204}; el nombre de subcasos limita la força de cada conclusió.}
\label{tab:experimental}
\end{table}

\begin{figure}[H]
\centering
\includegraphics[width=0.74\textwidth]{csmr_experimental_savings.png}
\caption{Estalvi dels presets experimentals. La mida de mostra s'indica sobre
cada barra i limita la força de la conclusió.}
\label{fig:experimental}
\end{figure}

\code{high\_ndvi} és el resultat més equilibrat: tres casos mid-roi i un
estalvi del 6,54\%. Per això es va promocionar a la configuració operativa del
benchmark Raspberry. \code{low\_ndvi} comprimeix més, però dos dels tres casos
tenen només 0,78\% de ROI; una regió tan dispersa és més vulnerable a
l'efecte dels blocs veïns. \code{dark\_regions} conserva ROI, però dedica
qualitat alta a gairebé dos terços de la imatge i deixa poc fons per degradar.

El cas \code{water\_body} requereix una precisió metodològica. Al full original
amb llindar 0,30, 119 dels 120 retalls no tenien ROI i un tenia una ROI de
0,098\%. L'optimització va reduir el llindar a 0,10, però només un crop
representatiu d'aigua havia estat seleccionat per aquell barratge i va passar
a CSMR. Per tant, una mostra no significa que la resta del dataset sigui tot
aigua o tota no-ROI amb el llindar nou; significa que l'experiment no va
regenerar els 120 retalls a 0,10. El 11,18\% és prometedor, però insuficient
per generalitzar.

\section{Preserve-roi}\label{sec:preserve_results}

La prova CSMR de \code{preserve-roi} va completar 78 subcasos. La
Taula~\ref{tab:preserve_cloud} mostra dos casos \code{cloud\_avoid} compartits
per les tres polítiques.

\begin{table}[H]
\centering
\small
\begin{tabular}{lrrr}
\toprule
Política & Estalvi vs \code{q204} & $\Delta$ROI vs \code{q204} &
$\Delta$fons vs \code{q204} \\
\midrule
\code{focus\_bgq128} & 17,06\% & +0,351 & -2,478 \\
\code{preserve Q240/Q128} & 15,88\% & +0,626 & -2,480 \\
\code{preserve Q255/Q128} & 15,30\% & +1,024 & -2,476 \\
\bottomrule
\end{tabular}
\caption{Compromís de \code{preserve-roi} en dos casos \code{cloud\_avoid}.
Els deltes són dB amb màscara comuna i el baseline és \code{q204}.}
\label{tab:preserve_cloud}
\end{table}

\begin{figure}[H]
\centering
\includegraphics[width=0.76\textwidth]{preserve_roi_tradeoff.png}
\caption{Intercanvi entre estalvi i preservació de ROI. La posició més a la
dreta afavoreix compressió; la posició més alta afavoreix qualitat ROI.}
\label{fig:preserve_tradeoff}
\end{figure}

El resultat respon a objectius diferents. \code{focus} és preferible si es
prioritza bitrate. Q255 és preferible si es vol assegurar una ROI clarament
millor que \code{q204}. Q240 és un compromís intermedi. No hi ha una política
dominant en les dues dimensions: la decisió depèn de si la missió valora més
baixar volum de dades o protegir la regió seleccionada.

\section{Evidència visual}\label{sec:visual_results}

Els panells següents tenen quatre columnes:

\begin{enumerate}[leftmargin=*]
  \item composició RGB de l'entrada;
  \item màscara del preset superposada en verd;
  \item Q-map ampliat, on colors clars indiquen Q més alta;
  \item error absolut amplificat per fer visible la seva distribució.
\end{enumerate}

L'error amplificat no és l'aspecte de la reconstrucció. La multiplicació visual
serveix per localitzar on es concentra la pèrdua. La imatge reconstruïda
completa acostuma a semblar gairebé igual que l'original; la prova quantitativa
és el PSNR i la prova espacial és el Q-map.

\begin{table}[H]
\centering
\scriptsize
\begin{tabular}{lrrrr}
\toprule
Panell visual & ROI \% & Estalvi vs \code{q204} & $\Delta$ROI &
$\Delta$fons \\
\midrule
\code{cloud\_avoid + focus} & 15,1 & 24,5\% & +0,474 & -2,660 \\
\code{local\_contrast + focus} & 10,1 & 20,7\% & +0,101 & -2,773 \\
\code{vegetation + focus} & 10,5 & 9,4\% & +0,156 & -1,921 \\
\code{high\_ndvi + focus} & 10,9 & 8,1\% & +0,161 & -2,297 \\
\code{cloud\_avoid + preserve Q240} & 15,1 & 22,9\% & +0,624 & -2,663 \\
\bottomrule
\end{tabular}
\caption{Lectura quantitativa dels panells visuals. El baseline és
\code{q204}; els deltes de PSNR són en dB amb la mateixa màscara ROI/fons que
mostra cada panel.}
\label{tab:visual_metrics}
\end{table}

\begin{figure}[H]
\centering
\includegraphics[width=\textwidth]{01_main_cloud_avoid_focus_bgq128.png}
\caption{\code{cloud\_avoid + focus\_bgq128}. El preset selecciona la ROI
detectada com a zona no ennuvolada o útil i el fons queda a Q128. En aquest
crop, la cobertura és 15,1\%, l'estalvi respecte de \code{q204} és 24,5\% i la
ROI millora 0,47~dB.}
\label{fig:visual_cloud}
\end{figure}

La Figura~\ref{fig:visual_cloud} és un cas especialment favorable, no la
mitjana del preset. La frontera es veu esglaonada perquè la decisió es pren en
blocs de $16\times16$. El Q-map mostra variació dins la ROI perquè \code{focus}
parteix d'\code{adaptive\_s8}; no és un nivell fix.

\begin{figure}[H]
\centering
\includegraphics[width=\textwidth]{02_main_local_contrast_focus_bgq128.png}
\caption{\code{local\_contrast + focus\_bgq128}. Les zones amb més variació
visible reben la qualitat adaptativa reforçada.}
\label{fig:visual_contrast}
\end{figure}

\code{local\_contrast} no etiqueta una coberta física. La màscara de la
Figura~\ref{fig:visual_contrast} indica on la desviació visible supera 0,035.
És útil per conservar vores i textures, però la missió hauria de decidir si
aquesta prioritat és coherent amb l'aplicació final.

\begin{figure}[H]
\centering
\includegraphics[width=\textwidth]{03_main_vegetation_focus_bgq128.png}
\caption{\code{vegetation + focus\_bgq128}. La màscara NDVI protegeix blocs
vegetats i deixa la resta a Q128.}
\label{fig:visual_vegetation}
\end{figure}

La Figura~\ref{fig:visual_vegetation} mostra una selecció més fragmentada.
Aquest tipus de geometria és rellevant perquè el camp receptiu travessa blocs
veïns. Els resultats indiquen ROI conservada en la mostra CSMR, però regions
encara més petites podrien requerir dilatar la màscara o utilitzar
\code{preserve-roi}.

\begin{figure}[H]
\centering
\includegraphics[width=\textwidth]{04_experimental_high_ndvi_focus_bgq128.png}
\caption{\code{high\_ndvi + focus\_bgq128}. Exemple del preset experimental que
va obtenir tres casos mid-roi consistents.}
\label{fig:visual_high_ndvi}
\end{figure}

\code{high\_ndvi} és més restrictiu que \code{vegetation}. La figura permet
entendre per què els dos presets poden coexistir: no són dues polítiques, sinó
dos nivells de selecció de la mateixa magnitud espectral.

\begin{figure}[H]
\centering
\includegraphics[width=\textwidth]{05_preserve_cloud_avoid_preserve_roi_q240_bgq128.png}
\caption{\code{cloud\_avoid + preserve-roi Q240/Q128}. La ROI té un color
uniforme perquè Q240 és fixa; aquesta és la diferència visual amb el Q-map
adaptatiu de la Figura~\ref{fig:visual_cloud}.}
\label{fig:visual_preserve}
\end{figure}

\section{Abast de la lectura}\label{sec:statistical_reading}

Les comparacions CSMR són aparellades per retall: cada resultat focus,
experimental o preserve-roi es compara amb el \code{q204} del mateix crop i la
mateixa màscara. Tot i això, les taules combinen poblacions diferents. El full
C-only conté 120 retalls; el CSMR principal en conté 20; les taules per preset
agreguen entre tres i cinc subcasos; i les auditories experimentals provenen de
mostres separades. Per això els percentatges s'han de llegir com a evidència de
mecanisme i de compromís, no com una mitjana universal de missió.

\section{Resultats que informen l'usuari}\label{sec:user_implications}

La geometria de la ROI canvia el valor del sistema. La
Taula~\ref{tab:user_roi_cases} resumeix les decisions derivades.

\begin{table}[H]
\centering
\small
\begin{tabularx}{\textwidth}{p{3.0cm}p{4.0cm}X}
\toprule
Cobertura observada & Conseqüència & Decisió recomanada \\
\midrule
ROI buida & Tot rep qualitat de fons; no es preserva cap regió. & Usar
\code{q204},
un altre preset o no transmetre si la missió ho permet. \\
ROI molt petita i dispersa & El camp receptiu travessa blocs veïns i pot
degradar la regió tot i el Q alt. & Dilatar màscara, usar preserve-roi o revisar
el llindar. \\
ROI intermèdia & Hi ha prou ROI i fons per redistribuir qualitat. & Focus és
el candidat principal si els deltes CSMR són favorables. \\
ROI molt ampla & Queda poc fons on estalviar i el bitrate s'apropa al control. &
Usar un llindar més selectiu o \code{q204}. \\
ROI total & Tots els blocs són prioritaris; no hi ha redistribució espacial. &
Usar qualitat global o preserve només si es vol Q alta uniforme. \\
\bottomrule
\end{tabularx}
\caption{Implicacions operatives de la cobertura produïda pel preset.}
\label{tab:user_roi_cases}
\end{table}

Aquestes regles mostren per què un preset no s'ha de jutjar només pel bitrate.
\code{low\_ndvi} aconsegueix estalvi en tres casos, però la ROI mitjana és
petita; \code{dark\_regions} conserva ROI, però n'ocupa massa; \code{water\_body}
té un cas prometedor sense cobertura suficient. El mode útil no és el que
produeix el fitxer més petit, sinó el que selecciona una regió coherent amb la
missió i la conserva a un cost assumible.

\subsection{Focus davant preserve-roi}

La comparació de la Taula~\ref{tab:preserve_cloud} permet formular una regla
d'usuari. Si el downlink és el recurs dominant, focus ofereix més estalvi. Si
la pèrdua d'una estructura ROI és més costosa que enviar alguns bits
addicionals, Q240 o Q255 proporcionen més marge. La selecció no s'ha
d'automatitzar només amb una mitjana global; forma part dels requisits de
missió.

Els resultats d'aquest capítol mostren el compromís entre qualitat espacial i
bitrate real. La pregunta següent és si aquest compromís té un cost assumible en
hardware de baix recurs; això es tracta al Capítol~\ref{ch:bord}.
